\section{Coordinate systems, trigonometry, and changes of variables}
\label{sec:coordinates}

This chapter addresses the Cartesian, polar, cylindrical and spherical
coordinate links of the source article. Its algebraic direction coordinates
are especially useful constructively: an angle may be represented by its
unit direction $(c,s)$ without choosing a discontinuous global argument.
We also construct the conventional sine and cosine functions with an
explicit finite-series bound, so that the forward radian formulas have
their usual meaning.

\subsection{Radian functions from bounded rational sums}

Identify a pair $(a,b)$ with the formal complex number $a+ib$ only as
notation for the rational multiplication rule
$(a,b)(c,d)=(ac-bd,ad+bc)$. Put
\[
 E_N(t)=\sum_{j=0}^{N}\frac{(it)^j}{j!}=C_N(t)+iS_N(t),
\]
where $C_N$ collects the even powers and $S_N$ the odd powers. These are
finite rational expressions for rational $t$. The norm used for estimates
in this subsection is $|(a,b)|_1=|a|+|b|$, which is submultiplicative by
the rational triangle inequality.

\begin{lemma}[A primitive-recursive trigonometric tail bound]
\label{coord:tail}
Let $M\ge1$ be an integer, $|t|\le M$, and $J\ge N\ge2M$. Then
\[
 |E_J(t)-E_N(t)|_1\le 2^{4M-N}.
\]
Consequently the schedule $N(M,n)=4M+n+6$ defines PR names
$C_{N(M,n)}(t)$ and $S_{N(M,n)}(t)$ for $\cos t$ and $\sin t$.
Different admissible choices of $M$ give equal names.
\end{lemma}
\begin{proof}
For every finite $K$, $j!\ge2^{j-1}$ for $j\ge1$, by induction on $j$.
Hence $\sum_{j=0}^K1/j!\le3$. The multinomial identity, obtained by
repeated distributivity, gives
\[
 \sum_{j=0}^{K}\frac{M^j}{j!}
 \le\left(\sum_{j=0}^{K}\frac1{j!}\right)^M
 \le3^M\le2^{2M}.
\]
The first inequality holds because all terms whose total degree is at
most $K$ occur in the product, with the displayed multinomial coefficients;
the additional terms are nonnegative. Starting at index $2M$, successive
terms $M^j/j!$ decrease by a factor at most $1/2$. Therefore
\[
 \sum_{j=N+1}^{J}\frac{M^j}{j!}
 \le 2^{2M}\sum_{j=N+1}^{J}2^{-(j-2M)}
 \le2^{4M-N}.
\]
All sums here are finite, and the geometric-sum estimate is an induction
identity. The pair triangle inequality proves the assertion for $E$.
The proposed schedule leaves error at most $2^{-n-6}$ after any further
truncation. Comparing two truncations at a common later degree proves
the name condition and independence of $M$. Powers, factorials, finite
sums and the degree schedule are primitive recursive. The multinomial
and geometric-sum proofs use quantifier-free induction on their finite
counters.
\end{proof}

\begin{theorem}[Addition laws and the unit-circle identity]
\label{coord:addition}
The constructed functions satisfy
\begin{align*}
 \cos(s+t)&=\cos s\cos t-\sin s\sin t,\\
 \sin(s+t)&=\sin s\cos t+\cos s\sin t,\\
 \cos^2t+\sin^2t&=1.
\end{align*}
Also $\cos0=1$, $\sin0=0$, $\cos(-t)=\cos t$, and
$\sin(-t)=-\sin t$. The assertions are arithmetic identities of fixed PR
names, first for rational parameters and then for fixed PR named inputs
with supplied magnitude bounds.
\end{theorem}
\begin{proof}
Multiply $E_N(s)E_N(t)$. For total degree $k\le N$, the coefficient is
\[
 i^k\sum_{j=0}^k\frac{s^jt^{k-j}}{j!(k-j)!}
       =i^k\frac{(s+t)^k}{k!}
\]
by the binomial formula. Thus its terms through total degree $N$ equal
$E_N(s+t)$. With an integer $M\ge|s|+|t|+1$, the sum of the absolute
values of the remaining coefficients is bounded by
$\sum_{k=N+1}^{2N}M^k/k!\le2^{4M-N}$.
Lemma~\ref{coord:tail} and the multiplication estimate of
Theorem~\ref{found:operations} give the addition laws at each requested
precision. The parity statements hold term by term in the finite sums.
Set $s=-t$ in the addition formula and use $E_N(0)=(1,0)$ to get the
unit-circle identity. This also proves $|\cos t|,|\sin t|\le1$ from
nonnegativity of squares, through finite rational observations.

To admit named arguments, note that for $|u|,|v|\le M$,
\[
 |E_N(u)-E_N(v)|_1
 \le |u-v|\sum_{j=1}^N\frac{jM^{j-1}}{j!}
 \le3^M|u-v|.
\]
The first inequality is the factorization of $u^j-v^j$; the second is
the finite bound in Lemma~\ref{coord:tail}. Query an argument at precision
$n+\ell(3^{M+1})+8$, allow the approximating value the bound $M+1$,
and use degree $4(M+1)+n+8$. These are PR schedules. The estimates
transfer the rational formulas with room for every observation error.
\end{proof}

\begin{theorem}[Quantitative trigonometric derivatives]
\label{coord:derivatives}
For $|h|\le1$,
\begin{align*}
 |\sin(t+h)-\sin t-h\cos t|&\le6h^2,\\
 |\cos(t+h)-\cos t+h\sin t|&\le6h^2.
\end{align*}
In particular the names have derivatives $\cos t$ and $-\sin t$ in the
finite-error sense, with an explicit modulus.
\end{theorem}
\begin{proof}
The finite series and Lemma~\ref{coord:tail} imply
$|\cos h-1|\le3h^2$ and $|\sin h-h|\le3|h|^3$ for $|h|\le1$:
factor $h^2$ or $|h|^3$ from the corresponding tail and bound the
remaining factorial sum by $3$. For a named $h$ with $|h|\le1$, clip its rational observations to $[-1,1]$ before applying these finite estimates; clipping preserves the name and the error bound. Substitute these estimates into the
addition formulas. Each coefficient $\sin t$ or $\cos t$ has absolute
value at most one. The first residual is
$\sin t(\cos h-1)+\cos t(\sin h-h)$, so it is at most
$3h^2+3|h|^3\le6h^2$; the second is handled by the displayed
cosine addition identity. For a difference quotient error at most
$2^{-n}$ it suffices to supply $0<|h|\le2^{-n}/6$; the quotient is
used only under this nonzero guard. The remainder inequalities themselves
include $h=0$ and avoid division entirely.
\end{proof}

\subsection{Polar and spatial coordinates}

\begin{theorem}[Polar direction coordinates and reconstruction]
\label{coord:polar}
For $r\ge0$ and $c^2+s^2=1$, put $(x,y)=(rc,rs)$. Then
$x^2+y^2=r^2$. Conversely, for a rational point $(x,y)\ne(0,0)$, the
constructed values
\[
 r=\sqrt{x^2+y^2},\qquad c=x/r,\qquad s=y/r
\]
give a unique radius and unit direction and reconstruct the point. At
the origin one may output $r=0,(c,s)=(1,0)$; the direction there is not
unique. For arbitrary fixed PR named points the same converse requires
a supplied positive lower bound on $x^2+y^2$, or the supplied assertion
that both coordinates are zero. The forward radian formulas are
$x=r\cos\theta$, $y=r\sin\theta$.
\end{theorem}
\begin{proof}
The forward assertion is the polynomial identity
$(rc)^2+(rs)^2=r^2(c^2+s^2)$. For a nonzero rational point the positive
rational $q=x^2+y^2$ has a computable rational lower bound. If
$q\ge\delta>0$, then $r\ge\min(1,\delta)$, since
$0<\delta\le1$ implies $\sqrt\delta\ge\delta$, and
$\delta\ge1$ implies $\sqrt\delta\ge1$. A dyadic lower bound below
$\min(1,\delta)$ is found by bounded search. Thus division by $r$ is
permitted by Theorem~\ref{found:operations}. Squaring proves the unit
identity, and multiplying gives reconstruction. Nonnegative-root
uniqueness forces any other nonnegative radius to be $r$, and the
nonzero radius then forces the same direction. The zero case is checked
directly. Theorem~\ref{coord:addition} supplies the forward radian case.
\end{proof}

The usual scalar formula $\theta=\arctan(y/x)$ requires both a chart and
a quadrant convention. No global scalar inverse is claimed here.
Direction coordinates encode the complete ray, including vertical and
negative-horizontal directions. The rational parametrizations of
Theorem~\ref{con:standard-param} give stereographic charts; their inverse
requires the chart denominator to be apart from zero. This is an atlas
interpretation, not a choice of a global continuous argument.

\begin{theorem}[Cylindrical coordinates]
\label{coord:cylindrical}
The transformation
\[
 (r,c,s,z)\longmapsto(rc,rs,z),\qquad r\ge0,\quad c^2+s^2=1,
\]
satisfies $x^2+y^2=r^2$. It has the inverse of
Theorem~\ref{coord:polar} on each horizontal plane, with the same axis
qualification. Hence $x^2+y^2=R^2$ is the circular cylinder of radius
$R\ge0$, and $z=k$ fixes a horizontal plane. Radian input substitutes
$c=\cos\theta,s=\sin\theta$.
\end{theorem}
\begin{proof}
Apply the polar theorem to the first two coordinates and leave $z$
unchanged. The inverse preserves that last coordinate. A fixed radius
condition places no restriction on $z$; fixing $z$ places no restriction
on the first two coordinates. These observations prove the locus
descriptions by the displayed equations. All inverse branches on
rational Cartesian inputs are rational zero tests on $x^2+y^2$.
\end{proof}

\begin{theorem}[Spherical coordinates with explicit axis data]
\label{coord:spherical}
For $\rho\ge0$, $c_\theta^2+s_\theta^2=1$, and
$c_\phi^2+s_\phi^2=1$ with $s_\phi\ge0$, set
\[
 x=\rho s_\phi c_\theta,\quad
 y=\rho s_\phi s_\theta,\quad z=\rho c_\phi.
\]
Then $x^2+y^2+z^2=\rho^2$ and the cylindrical radius is
$r=\rho s_\phi$. For a rational point away from the origin and the
vertical axis an inverse is
\[
 \rho=\sqrt{x^2+y^2+z^2},\quad r=\sqrt{x^2+y^2},\quad
 (c_\theta,s_\theta)=(x/r,y/r),\quad
 (c_\phi,s_\phi)=(z/\rho,r/\rho).
\]
On the nonzero vertical axis take $(c_\theta,s_\theta)=(1,0)$,
$s_\phi=0,c_\phi=\sgn z$. At the origin use $\rho=0$ and arbitrary
fixed admissible directions. The inverse for general named data
requires the corresponding radial gaps or certified axis cases.
\end{theorem}
\begin{proof}
Substitution and the two unit identities give
\[
 x^2+y^2+z^2
 =\rho^2s_\phi^2(c_\theta^2+s_\theta^2)+\rho^2c_\phi^2
 =\rho^2.
\]
The horizontal norm squared is $\rho^2s_\phi^2$; since both factors
are nonnegative, its nonnegative root is $\rho s_\phi$. The inverse
unit identities follow from $r^2=x^2+y^2$ and $\rho^2=r^2+z^2$;
substitution then recovers all three coordinates. The finitely many
axis cases are explicit rational branches. The nonzero denominators
receive gaps exactly as in the polar theorem. For conventional
inclination and azimuth one substitutes
$c_\phi=\cos\phi,s_\phi=\sin\phi$ and
$c_\theta=\cos\theta,s_\theta=\sin\theta$, retaining the
nonnegative-inclination-sine condition.
\end{proof}

\subsection{Oblique frames and transformation of equations}

\begin{theorem}[The metric in an oblique coordinate frame]
\label{coord:gram-metric}
Let $e_1,\ldots,e_d$ be an invertible rational frame, $d=2$ or $3$, with
matrix $M$ having those columns. Coordinates $u$ represent $P=O+Mu$.
The squared distance between coordinates $u,v$ is
\[
 (u-v)^TG(u-v),\qquad G=M^TM.
\]
The form $G$ is positive definite. It is the identity precisely when
the frame is orthonormal. Changes of origin leave $G$ unchanged.
\end{theorem}
\begin{proof}
The physical displacement is $M(u-v)$, whose norm square is
$(u-v)^TM^TM(u-v)$. If this sum of squares is zero, every coordinate
of $M(u-v)$ is zero; multiplication by the rational inverse of $M$
gives $u=v$. Thus the form is positive on nonzero rational vectors;
for names a quantitative nonzero vector is needed to infer a positive
gap, while zero-kernel injectivity remains valid. The $(i,j)$ entry of
$G$ is $e_i\cdot e_j$, so $G=I$ is exactly the orthonormality
condition. The origin cancels in a displacement. This explains why a
Euclidean distance formula cannot be reused unchanged in an arbitrary
affine frame.
\end{proof}

\begin{theorem}[Transforming a locus uses the inverse substitution]
\label{coord:locus-transform}
Let $T(u)=Mu+t$ be an invertible rational affine transformation and let
a locus be specified by $F(u)=0$. Then its transformed locus is
\[
 F\bigl(M^{-1}(x-t)\bigr)=0.
\]
In particular translation by $(h,k)$ gives $F(x-h,y-k)=0$;
scaling by nonzero $(b,a)$ gives $F(x/b,y/a)=0$; rotation with
unit pair $(c,s)$ gives $F(cx+sy,-sx+cy)=0$.
\end{theorem}
\begin{proof}
If $u$ is on the original locus and $x=T(u)$, then
$M^{-1}(x-t)=u$, which proves the transformed equation. Conversely,
that equation supplies the explicit preimage $u=M^{-1}(x-t)$, so
the point is an image of a point of the locus. The displayed special
cases insert the inverses of a translation, a diagonal matrix, and
$\left(\begin{smallmatrix}c&-s\\s&c\end{smallmatrix}\right)$,
respectively. The latter inverse is its transpose because $c^2+s^2=1$.
For a graph $y=f(x)$ this also proves that
$y=a f(b(x-k))+h$ corresponds to the point map
$(u,v)\mapsto(k+u/b,h+av)$, under $a,b\ne0$. If $f$ is a
fixed represented function with primitive-recursive extensional evaluation
(equal input names give equal output names), the assertion is a substitution
schema for that function. No universal evaluator for all functions is used.
\end{proof}

\begin{proposition}[Coordinate Jacobian identities]
\label{coord:jacobians}
With the conventional radian formulas, the polar Jacobian determinant
for input order $(r,\theta)$ is $r$, the cylindrical determinant for
$(r,\theta,z)$ is $r$, and the spherical determinant for
$(\rho,\phi,\theta)$ is $\rho^2\sin\phi$. Their absolute values
give the associated local area or volume scaling. These assertions
are determinant identities; no integration change-of-variables theorem
is assumed.
\end{proposition}
\begin{proof}
The polar derivative columns are $(\cos\theta,\sin\theta)$ and
$(-r\sin\theta,r\cos\theta)$, whose determinant is
$r(\cos^2\theta+\sin^2\theta)=r$. The cylindrical derivative
adds the independent column $(0,0,1)$, leaving the same determinant.
For spherical coordinates put $s=\sin\phi,c=\cos\phi$ and
$C=\cos\theta,S=\sin\theta$. The three columns are
\[
 (sC,sS,c),\quad(\rho cC,\rho cS,-\rho s),\quad
 (-\rho sS,\rho sC,0).
\]
The cross product of the last two is
$\rho^2s(sC,sS,c)$. Its dot product with the first is
$\rho^2s(s^2+c^2)=\rho^2s$. The derivative formulas follow from
Theorem~\ref{coord:derivatives} and finite products. Swapping the
order of the two angle variables changes the determinant's sign.
At $r=0$, or $\rho\sin\phi=0$, these formulas explicitly expose
the chart singularity. Local parallelepiped volume scaling is the
determinant identity of Theorem~\ref{geo:cross}.
\end{proof}

The proofs in this chapter use finite sums, factorials, powers,
polynomial identities, explicit root brackets and guarded reciprocals.
The only variable-length proofs are the finite sum estimates, governed
by Lemma~\ref{found:loops}. They therefore respect the same first-order
\ipra\ discipline as the exact geometry chapters.
